\newcommand{\verfasser}{Phil Steinhorst}					% Verfasser des Dokuments
\newcommand{\fach}{Ein Skript}								% Titel der Vorlesung								
\newcommand{\shortFach}{Skript}								% Titel kurz
\newcommand{\prof}{Prof.\,Dr.\,Schlotze}					% Dozent
\newcommand{\untertitel}{gelesen von \prof}					% Untertitel
\newcommand{\semester}{Sommersemester 2014}					% Semester
\newcommand{\homepage}{http://keen.bplaced.net/augenkrebs/augenkrebs.htm}	% Vorlesungshomepage
% % % % % % % % % % % % % % % % % % % % % % % % % % % % % % % % %
%		PhistScript.tex											%
%		XeLaTeX-Settings for Scripts							%
%																%
%		Author: Phil Steinhorst									%
% % % % % % % % % % % % % % % % % % % % % % % % % % % % % % % % %

\setlength\parindent{0pt}             % ohne Einrueckung

% % % Konfiguration von scrheadings
\setheadsepline{1pt}[\color{black}]
\pagestyle{scrheadings}
\clearscrheadfoot

% % % Kopf-/Fußzeilenlayout
\providecommand{\shortFach}{\fach}
\ohead{\small \leftmark}
\ofoot[{ \small \thepage}]{ \small \thepage} 
\automark{section}

% Metadaten
\title{\fach}
\author{\verfasser}
\subtitle{\untertitel}
\date{\semester}

% Inhaltsverzeichnis
\setcounter{tocdepth}{1}
\usepackage[tocindentauto]{tocstyle}
\usetocstyle{KOMAlike}	

% Nummerierung
\numberwithin{equation}{section}		% Section in Equation-Nummerierung einbeziehen

% ntheorem package
\newcounter{tcount}						% Zähler für Theoreme
\numberwithin{tcount}{section}			% Section in Theorem-Nummerierung einbeziehen
\usepackage[hyperref]{ntheorem}			% Lade ntheorem mit hyperref-Workaround
\theoremstyle{break}					% Zeilenumbruch nach Theoremkopf
\theorembodyfont{\normalfont}			% nicht kursiv
\theorempreskipamount0.5cm				% Abstand vor Theorem
\theorempostskipamount0.5cm				% Abstand nach Theorem
\newtheorem{thm}[tcount]{Theorem}
\newtheorem{bsp}[tcount]{Beispiel}
\newtheorem{defn}[tcount]{Definition}
\newtheorem{bem}[tcount]{Bemerkung}
\newtheorem{lemma}[tcount]{Lemma}

% Workaround für Tabulatoren im Theorem-Verzeichnis
\makeatletter
\def\thm@@thmline@name#1#2#3#4{%
        \@dottedtocline{-2}{0em}{2.3em}%
                   {\makebox[\widesttheorem][l]{#1 \protect\numberline{#2}}#3}%
                   {#4}}
\@ifpackageloaded{hyperref}{
\def\thm@@thmline@name#1#2#3#4#5{%
    \ifx\\#5\\%
        \@dottedtocline{-2}{0em}{2.3em}%
            {\makebox[\widesttheorem][l]{#1 \protect\numberline{#2}}#3}%
            {#4}
    \else
        \ifHy@linktocpage\relax\relax
            \@dottedtocline{-2}{0em}{2.3em}%
                {\makebox[\widesttheorem][l]{#1 \protect\numberline{#2}}#3}%
                {\hyper@linkstart{link}{#5}{#4}\hyper@linkend}%
        \else
            \@dottedtocline{-2}{0em}{2.3em}%
                {\hyper@linkstart{link}{#5}%
                  {\makebox[\widesttheorem][l]{#1 \protect\numberline{#2}}#3}\hyper@linkend}%
                    {#4}%
        \fi
    \fi}
}
\makeatother
\newlength\widesttheorem
\AtBeginDocument{
  \settowidth{\widesttheorem}{Proposition 10.10\quad}	% Referenzstring für die Breite der ersten Spalte
}

\theoremlisttype{optname}	% nur Theoreme auflisten, die einen Namen haben. 							% XeLaTeX-Konfiguration für Skripte

% % % Dokumentspezifische Einstellungen % % %
\setcounter{tocdepth}{2}				% Tiefe im Inhaltsverzeichnis (1: nur Sections, 2: auch subsections...)
\numberwithin{equation}{section}		% Section in Equation-Nummerierung einbeziehen
\newcounter{tcount}						% Zähler für Theoreme
\numberwithin{tcount}{section}			% Section in Theorem-Nummerierung einbeziehen

% ab hier können neue Umgebungen für Sätze definiert und benannt werden!
% Syntax: \newtheorem{interner Umgebungsname}[Zähler]{Name im Dokument}
% Möchte man Sätze und Definitionen getrennt durchnummerieren, benötigt man weitere Zähler.
\newtheorem{satz}[tcount]{Satz}
\newtheorem{bsp}[tcount]{Beispiel}
% % % % % % % % % % % % % % % % % % % % % % %

\begin{document}
% % % % % % % % % % % % % % % % % % % % % % % % % % % % % % % % %
%		PhistTitle.tex											%
%		Layout für Skript-Titelseiten							%
%																%
%		Author: Phil Steinhorst									%
% % % % % % % % % % % % % % % % % % % % % % % % % % % % % % % % %
\begin{titlepage}
\pagestyle{empty}
\begin{center}
\begin{minipage}{0.4\textwidth}
\begin{flushleft}
\includegraphics[height=1.5cm,keepaspectratio]{../!config/wwulogo.pdf}
\end{flushleft}
\end{minipage}
\hfill
\begin{minipage}{0.4\textwidth}
\begin{flushright}
\vspace*{0.3cm}
\includegraphics[height=1.2cm,keepaspectratio]{../!config/fb10logo.pdf} \
\end{flushright}
\end{minipage}

% Title
\vspace*{6cm}
\textbf{\Huge{\fach}} \\
\vspace{0.2cm} 
\textbf{{\LARGE \untertitel}} \\
\vspace{0.6cm}
\LARGE{Mitschrift von \verfasser} \\
\vspace{0.6cm}
\LARGE{\semester} \\
\vspace*{2cm}
\footnotesize{\url{\homepage}}
\vfill

\vspace*{2cm}
\begin{flushright}
{\footnotesize Stand: \today}
\end{flushright}
\end{center}
\end{titlepage}							% Titelseite ausgeben
\section*{Vorwort}
	Das ist ein Vorwort.
\newpage
	
\tableofcontents											% Inhaltsverzeichnis ausgeben
\newpage

\section{Ein Abschnitt}
	Hier beginnt ein Abschnitt.
	
\subsection{Ein Unterabschnitt}
	... und hier ein Unterabschnitt.
	
\begin{satz}[Ein wichtiger Satz]
	Blaukraut bleibt Blaukraut und \Index{Brautkleid} bleibt Brautkleid. \index{Blaukraut}
\end{satz}

\newpage
\appendix							% Anhang beginnt hier (Nummerierung mit Buchstaben)

\section{Ein Anhang}
	Ich bin der Anhang.

\printindex							% Index ausgeben
\subsection*{Liste der Sätze und Definitionen}
\listtheorems{satz}					% listet angegebene Sätze mit Namen und Seitenzahl auf
\end{document}